%% Fragments — task proposal for the ELOQUENT lab, CLEF 2027
%% Structured after the CLEF call for lab proposals (3–4 pages, eleven required items),
%% restricted to the items that concern a single task within an existing lab.
\documentclass[10pt,a4paper]{article}
\usepackage[margin=1.7cm]{geometry}
\usepackage[T1]{fontenc}
\usepackage[utf8]{inputenc}
\usepackage{lmodern}
\usepackage{microtype}
\usepackage{booktabs}
\usepackage{enumitem}
\usepackage{amsmath,amssymb}
\usepackage{listings}
\usepackage{xcolor}
\usepackage[hidelinks]{hyperref}
\usepackage[round]{natbib}
\setlist{nosep,leftmargin=1.4em}
\lstset{basicstyle=\footnotesize\ttfamily,frame=single,xleftmargin=0.6em,aboveskip=4pt,belowskip=4pt,breaklines=true,columns=fullflexible,keepspaces=true}
\setlength{\parskip}{2pt plus 1pt}
\usepackage{titlesec}\titlespacing*{\section}{0pt}{7pt}{3pt}\titlespacing*{\subsection}{0pt}{5pt}{2pt}
\setlength{\parindent}{0pt}

\title{\textbf{Fragments}: Can a Language Model Learn a Language from Its Specification Alone?\\[3pt]
\large A task proposal for the ELOQUENT lab at CLEF 2027}
\author{Stergios Chatzikyriakidis\thanks{University of Crete, \texttt{stergios.chatzikyriakidis@uoc.gr}}
\and Aarne Talman\thanks{University of Helsinki}}
\date{\vspace{-4pt}Draft of 29 September 2026}

\begin{document}
\maketitle

\section{Summary and relevance to ELOQUENT}
ELOQUENT evaluates high-level quality criteria of generative language models through open-ended shared tasks with light-weight participation and automatic or semi-automatic scoring \citep{karlgren2024eloquent,karlgren2025eloquent}. Fragments adds a criterion no current task covers: \textbf{specification-following}. At test time the organisers release a new formal \emph{fragment} of a language in Montague's sense \citep{montague1973}: a small grammar, a lexicon of invented words, compositional semantic rules from syntax to first-order logic, and a handful of finite models. Systems must parse sentences of the fragment, produce their meanings, and decide entailments with a justification a machine can check. Because the fragment is new each year it cannot be in any training set, and because every answer is checked by a parser, a model checker or a theorem prover, scoring needs no annotators. The task operationalises the question at the centre of the debate on the limits of language models, rule use versus pattern matching, with a score that is a proof rather than a judgement. It fits ELOQUENT's format exactly: JSONL in, scores out, an experiment report, and a builder/breaker option modelled on Voight-Kampff. It is not a competition: the interesting output is the map of \emph{which kinds of computation} recover a specification, so leaderboards are reported per system type.

\section{Motivation and novelty}
Existing evaluations test whether a model has absorbed a language's regularities from data. COGS and SLOG train the grammar in and test structural generalisation \citep{kim2020cogs,li2023slog}; MTOB gives a model an informal grammar book of Kalamang and scores translation \citep{tanzer2024mtob}; FraCaS and FOLIO are static, English, and label-scored \citep{cooper1996fracas,han2024folio}; LingOly-style puzzles test rule induction from a few examples \citep{bean2024lingoly}. In all of them the grammar, the words or the logic were in the training data, so a correct answer cannot distinguish computation from retrieval. Fragments removes the confound by construction: the rules are stated explicitly, the words are invented, and a \emph{counter-English} variant makes the familiar answer wrong. In the mathematical world the analogous move has been productive: miniF2F and ProofNet score autoformalisation by a proof kernel \citep{zheng2022minif2f,azerbayev2023proofnet}, but against a background library the model has seen; Fragments is their natural-language counterpart with an unseen background theory. The distinct claim is therefore \emph{in-context formal specification plus verified meaning and proof-scored inference}.

\section{Task definition}
\subsection{The fragment package}
Participants receive four plain-text files. A linguist reads them in ten minutes; that is the point. The toy fragment below is the one we would ship as the practice example.

\begin{lstlisting}
lexicon.txt    PN: Tam, Ren   N: blicket, wug   IV: daxes   TV: fepps   Det: kro, sib, nul
grammar.txt    S -> NP VP     NP -> Det N | PN     VP -> IV | TV NP
semantics.txt  [S NP VP] = NP(VP)      [NP Det N] = Det(N)      [VP TV NP] = \x. NP(\y. TV(x,y))
               kro = \P.\Q. forall x (P(x) -> Q(x))      sib = \P.\Q. exists x (P(x) & Q(x))
               nul = \P.\Q. ~exists x (P(x) & Q(x))      Tam = \P. P(tam)     daxes = \x. daxes(x)
models.json    M1: D={a,b,c} blicket={a,b} wug={c} daxes={a} fepps={(a,c),(b,c)} tam=a ren=c
               M2: D={a,b}   blicket={a,b} wug={}  daxes={a,b} fepps={}          tam=a ren=b
\end{lstlisting}

\subsection{Subtasks}
Items and answers are JSON lines; logical forms use TPTP first-order syntax, so an entry needs nothing but a model's outputs.
\begin{enumerate}
\item \textbf{Parse.} Decide whether a string belongs to the fragment and give its derivation. \texttt{kro blicket fepps sib wug} is in, with tree \texttt{[S [NP [Det kro] [N blicket]] [VP [TV fepps] [NP [Det sib] [N wug]]]]}; \texttt{blicket kro daxes} is out.
\item \textbf{Interpret.} Give the logical form: \texttt{!\,[X]\,:\,(blicket(X) => ?\,[Y]\,:\,(wug(Y) \& fepps(X,Y)))}. Any equivalent formula counts; nothing is string-matched.
\item \textbf{Infer.} Given premises \texttt{kro blicket fepps Ren}, \texttt{sib blicket daxes} and conclusion \texttt{sib blicket fepps Ren}, give the label (entailment) and a machine-checkable justification. In the \emph{tools track} the justification is the set of formulas, which the organisers run through a prover; in the \emph{no-tools track} it is a proof in a fixed ten-rule natural-deduction calculus.
\item \textbf{Break} (optional; builder/breaker as in Voight-Kampff). Submit new items with gold; an item is admitted only if the organisers' checker verifies its gold against the specification, and admitted items are given to the other systems.
\end{enumerate}

\subsection{Tracks and system types}
Two tracks in the first edition: \emph{no-tools} (a model's outputs from the specification in its context) and \emph{tools} (a model with a prover, or with code it writes). A \emph{constructive} track, where the semantics is dependent type theory in the tradition of \citet{ranta1994} and \citet{chatzikyriakidis2020} and submissions are proof terms checked by Coq or Lean, is planned for the second edition. Every entry declares its type: prompted model, prover loop, program synthesis (the model writes a parser and interpreter from the spec), neuro-symbolic split, or fine-tuned. Results are reported per type, because ``a model that writes a Python interpreter'' and ``a model that answers from the spec'' are both legitimate and both interesting, and averaging them would hide the finding. Two reference systems appear in every table: the \emph{compiled specification}, which scores 100\% by construction and shows the task is solvable from the spec alone, and a group of linguistics students, the human calibration for the difficulty ladder.

\section{Data}
\paragraph{Generation.} Items come from a generator, not from annotation, which is what makes the task cheap to repeat annually. Each cycle: (1)~sample a fragment (determiner meanings, verb arities, word order, nonce lexicon); (2)~generate sentences with their derivations and gold logical forms by compositional evaluation; (3)~sample finite models, some released and some hidden, and evaluate every gold formula on them; (4)~build inference items by sampling premise sets and conclusions and labelling them with Prover9 and Mace4 \citep{mccune2010prover9}, balancing the three labels and discarding items on which the prover does not terminate; (5)~produce the variants. The generator and the seed are published after each cycle.

\paragraph{Controlled variants.} Every structure appears with nonce words; with real words in English, Finnish, Czech and Greek, the same abstract grammar linearised through the Grammatical Framework and its Resource Grammar Library \citep{ranta2004gf}; with scrambled word order; and in a \emph{counter-English} version in which one semantic rule contradicts the familiar reading (\texttt{kro} means \emph{no}). Consistency across variants is the measurement, and the counter-English drop is the headline number per system. One control fragment per cycle has a deliberately inconsistent specification, to detect systems that ignore the spec.

\paragraph{Phenomena ladder.} Difficulty climbs from quantification (level~1) through negation and monotonicity (2), relative clauses and scope ambiguity, where gold is a set of readings (3), pronouns and donkey sentences with a dynamic semantics after \citet{groenendijk1991} (4), plurals and comparatives (5), to intensional contexts with worlds in the models (6). FraCaS sections give the template, and results are reported per level.

\paragraph{Gold and its verification.} Gold is never annotated; it is computed from the specification. The specification itself is authored as a Coq file (grammar as an inductive type, semantics as the interpretation function, models as instances) and exported to the four text files, so every gold logical form is a theorem and every gold label a prover verdict. The organisers' library of machine-checked semantic theories (FORMAL-ATLAS: Montague, Lambek, DPL, MTT, Ranta, inquisitive, probabilistic and categorical semantics; about 1,000 theorems, none admitted) is the source of the ladder's levels.

\section{Evaluation setup and metrics}
\begin{center}\small
\begin{tabular}{@{}p{2.6cm}p{6.2cm}p{6.2cm}@{}}
\toprule
Subtask & Checker & Metric \\
\midrule
Parse & the organisers' parser for the released grammar & grammaticality accuracy; derivation exact match \\
Interpret & model checking on released and hidden models; Prover9 on gold~$\leftrightarrow$~submitted as a bonus & truth-value agreement on every model; for ambiguous items, set equality of readings modulo equivalence, plus reading recall \\
Infer (tools) & Prover9 and Mace4 on the submitted formulas & label accuracy \\
Infer (no-tools) & proof checker for the fixed calculus & label accuracy and proof validity, reported separately \\
Break & the same checkers on the submitted gold & admitted items; fooling rate on other systems \\
\bottomrule
\end{tabular}
\end{center}
Three choices behind the table. Equivalence of first-order formulas is undecidable, so the primary Interpret score is semantic (truth on models), with hidden models so that a system cannot fit the visible situations. Ambiguous items carry a set of golds. The gap between label accuracy and proof validity in Infer separates reasoning from guessing and is the lab's central number. All scoring is deterministic and reproducible from the published generator.

\section{Target audience, outreach and expected participation}
The task serves model developers (a contamination-free measure of in-context rule-following with per-phenomenon diagnostics), linguists (an empirical arbiter with a proof kernel behind it), and the mathematical-AI community (an autoformalisation benchmark with an unseen background theory). The entry barrier is deliberately minimal: a prompting baseline (specification in context, few-shot from the practice fragment) and the scoring scripts are shipped with the practice fragment, so a team only needs to run a model and submit JSONL. We expect 8--15 teams: ELOQUENT's existing participants, the NLI and semantic-parsing communities (COGS, FraCaS, FOLIO users), the natural-logic and theorem-proving-for-NLI community (NALOMA), and mathematical-AI groups with prover-in-the-loop systems, for whom the tools track is a drop-in target. Outreach through the ELOQUENT channels, the NALOMA and IWCS communities, the Grammatical Framework community, and the organisers' networks in formal semantics.

\section{Organisation, timeline and resources}
\begin{center}\small
\begin{tabular}{@{}ll@{}}
\toprule
Autumn 2026 & task description to the ELOQUENT organisers; fragment generator written \\
January 2027 & practice fragment with full gold, scoring scripts, prompting baseline released \\
April 2027 & test fragment released; two-week submission window; Break submissions \\
May 2027 & scores returned; experiment reports due \\
Summer 2027 & task overview paper; working notes \\
14--17 September 2027 & CLEF 2027, Bucharest: results and the builder/breaker round \\
\bottomrule
\end{tabular}
\end{center}
Existing infrastructure: the NESTOR pipeline (natural language to FOL and Coq, Prover9/Mace4 scoring, four experimental conditions, a student evaluation protocol, containerised), built in the first organiser's seminar on computational semantics in Athens; a deployed Coq checker service; the FORMAL-ATLAS library; GF and the Resource Grammar Library for the multilingual variants. The one new component is the fragment generator (about a week of work) and, if the no-tools track runs in the first edition, the ten-rule proof checker. Compute for scoring is negligible: provers with timeouts on a few thousand items. A minimal viable first edition is Parse, Interpret and the tools-track Infer on the Montague fragment, which needs nothing beyond the generator.

\section{Organisers and prior experience}
\textbf{Stergios Chatzikyriakidis} (University of Crete): type-theoretic semantics, interactive theorem proving for natural language inference \citep{chatzikyriakidis2014,chatzikyriakidis2020}, and neurosymbolic NLP; organiser of CoNLI 2017, IWCS 2019 (Program and Local Chair), the CSTFRS series (IWCS 2019, ESSLLI 2021), NALOMA III (2022) and the fourth Natural Logic Meets Machine Learning workshop (2023); leads the NESTOR seminar project and FORMAL-ATLAS. \textbf{Aarne Talman}: co-founder of ELOQUENT and organiser of its first edition \citep{karlgren2024eloquent}, with a research record on the generalisation and robustness of natural language inference systems across benchmarks \citep{talman2019generalization}; brings the lab's own experience of running open-ended, automatically scored tasks and the link to the Robustness line. The students of the Athens seminar on computational semantics in which NESTOR was built will run the human-baseline study and assist with the practice release.

\section{New activity, relation to other efforts, and risks}
Fragments is a new task, not a continuation. It is complementary to Voight-Kampff (from which it borrows the builder/breaker design) and to the Robustness task (consistency across paraphrase; here consistency across variants with the semantics provably held fixed). It is also complementary to a workshop proposal by the first organiser (LingEmbed, EACL 2027), which asks whether a linguistic theory compiled \emph{inside} a model helps; Fragments hands the theory to an unmodified model as input and measures the in-context ceiling such a module has to beat, and never scores one theory against another.

Risks and mitigations. ``It is a synthetic puzzle'': the natural-language twins and the FraCaS-derived ladder tie every item to a documented linguistic phenomenon, and the human baseline shows it is learnable from the spec. ``Too high an entry barrier'': one fixed output format and a shipped baseline. ``Trust in the gold'': derived from a Coq specification and reproducible from the published generator. ``Too ambitious for a first edition'': the minimal viable task above, or as fallback proof-checked NLI alone (FraCaS-style items where systems must supply a verifiable derivation), which the existing pipeline already runs.

{\small\setlength{\bibsep}{1pt}
\bibliographystyle{plainnat}
\bibliography{refs}}
\end{document}
